% GENERATED FILE. Edit canonical lessons, then run npm run generate:books.
% canonical-source-hash: fnv1a64:50178dc5d5312f61
% canonical-lessons: BN-C75-dhoni, BN-C75-gorib, BN-C75-nongra, BN-C75-teto, BN-C75-tok

\chapter{Rich, Poor, Dirty, Bitter, Sour}
\label{ch:bn-rich-poor-dirty-bitter-sour}
\hlchaptermodality{75}

\begin{chapteropening}
\textbf{By the end of this chapter:} \emph{I can say rich, poor, dirty, bitter and sour.}

Complete the last lesson of chapter 75: I can say rich, poor, dirty, bitter and sour.
\end{chapteropening}

\section[\texorpdfstring{\bn{ধনী} (dhôni)}{ধনী (dhôni)}]{\bn{ধনী} (dhôni) — rich}

\label{lesson:BN-C75-dhoni}



\begin{quote}
Before the new one: say the Bengali for full, then the Bengali for empty.
\end{quote}

\subsection*{You'll want to know: \bn{ধনী}}
\textbf{\bn{ধনী}} — \emph{dhôni} — \textquotedblleft{}rich\textquotedblright{}.

\begin{cousinweb}
\textbf{\bn{ধনী}} (\emph{dhôni}) is inherited from Sanskrit \emph{dhanin}, \textquotedblleft{}rich\textquotedblright{}, and has come down with almost no change (CDIAL 6722).
\end{cousinweb}

\begin{grammarlens}[title={What you've built}]
One more word for this chapter.
\end{grammarlens}

\subsection*{Your turn}
\begin{itemize}
\raggedright
  \item \emph{Say it:} \emph{dhôni}
  \item \emph{Say it:} \emph{dhôni}, once more
  \item \emph{Say it:} say \emph{dhôni}
  \item \emph{Recall:} say the Bengali for full, then the Bengali for empty, then say \emph{dhôni} again
\end{itemize}

\begin{tcolorbox}[breakable,colback=teal!4,colframe=teal!35!black,title={Before you move on}]
What does \bn{ধনী} mean? (\textquotedblleft{}Rich\textquotedblright{}.) Say it once more.
\end{tcolorbox}
\section[\texorpdfstring{\bn{গরিব} (gôrib)}{গরিব (gôrib)}]{\bn{গরিব} (gôrib) — poor}

\label{lesson:BN-C75-gorib}



\begin{quote}
Before the new one: say the Bengali for empty, then the Bengali for rich.
\end{quote}

\subsection*{You'll want to know: \bn{গরিব}}
\textbf{\bn{গরিব}} — \emph{gôrib} — \textquotedblleft{}poor\textquotedblright{}.

\begin{cousinweb}
\textbf{\bn{গরিব}} (\emph{gôrib}) is a loan from Arabic \emph{gharīb}, \textquotedblleft{}a stranger\textquotedblright{}, which reached Bengali through Persian. Its Arabic root is \emph{gh-r-b} (Lane). In Bengali, as in Hindi, it has come to mean \textquotedblleft{}poor\textquotedblright{}.
\end{cousinweb}

\begin{grammarlens}[title={What you've built}]
One more word for this chapter.
\end{grammarlens}

\subsection*{Your turn}
\begin{itemize}
\raggedright
  \item \emph{Say it:} \emph{gôrib}
  \item \emph{Say it:} \emph{gôrib}, once more
  \item \emph{Say it:} say \emph{gôrib}
  \item \emph{Recall:} say the Bengali for empty, then the Bengali for rich, then say \emph{gôrib} again
\end{itemize}

\begin{tcolorbox}[breakable,colback=teal!4,colframe=teal!35!black,title={Before you move on}]
What does \bn{গরিব} mean? (\textquotedblleft{}Poor\textquotedblright{}.) Say it once more.
\end{tcolorbox}
\section[\texorpdfstring{\bn{নোংরা} (nongrā)}{নোংরা (nongrā)}]{\bn{নোংরা} (nongrā) — dirty}

\label{lesson:BN-C75-nongra}



\begin{quote}
Before the new one: say the Bengali for rich, then the Bengali for poor.
\end{quote}

\subsection*{You'll want to know: \bn{নোংরা}}
\textbf{\bn{নোংরা}} — \emph{nongrā} — \textquotedblleft{}dirty\textquotedblright{}.

\begin{cousinweb}
CDIAL lists Bengali \emph{noṅarā}, \textquotedblleft{}filthy\textquotedblright{}, under \emph{*luṅgara}, one of a group of forms for \textquotedblleft{}defective\textquotedblright{} that it reconstructs (CDIAL 11072).
\end{cousinweb}

\begin{grammarlens}[title={What you've built}]
One more word for this chapter.
\end{grammarlens}

\subsection*{Your turn}
\begin{itemize}
\raggedright
  \item \emph{Say it:} \emph{nongrā}
  \item \emph{Say it:} \emph{nongrā}, once more
  \item \emph{Say it:} say \emph{nongrā}
  \item \emph{Recall:} say the Bengali for rich, then the Bengali for poor, then say \emph{nongrā} again
\end{itemize}

\begin{tcolorbox}[breakable,colback=teal!4,colframe=teal!35!black,title={Before you move on}]
What does \bn{নোংরা} mean? (\textquotedblleft{}Dirty\textquotedblright{}.) Say it once more.
\end{tcolorbox}
\section[\texorpdfstring{\bn{তেতো} (teto)}{তেতো (teto)}]{\bn{তেতো} (teto) — bitter}

\label{lesson:BN-C75-teto}



\begin{quote}
Before the new one: say the Bengali for poor, then the Bengali for dirty.
\end{quote}

\subsection*{You'll want to know: \bn{তেতো}}
\textbf{\bn{তেতো}} — \emph{teto} — \textquotedblleft{}bitter\textquotedblright{}.

\begin{cousinweb}
\textbf{\bn{তেতো}} (\emph{teto}) is a tadbhava, a word inherited from Sanskrit and reshaped by sound change. It comes from \emph{tikta}, \textquotedblleft{}pungent, bitter\textquotedblright{} (CDIAL 5806). CDIAL writes the Bengali form \emph{tita}.
\end{cousinweb}

\begin{grammarlens}[title={What you've built}]
One more word for this chapter.
\end{grammarlens}

\subsection*{Your turn}
\begin{itemize}
\raggedright
  \item \emph{Say it:} \emph{teto}
  \item \emph{Say it:} \emph{teto}, once more
  \item \emph{Say it:} say \emph{teto}
  \item \emph{Recall:} say the Bengali for poor, then the Bengali for dirty, then say \emph{teto} again
\end{itemize}

\begin{tcolorbox}[breakable,colback=teal!4,colframe=teal!35!black,title={Before you move on}]
What does \bn{তেতো} mean? (\textquotedblleft{}Bitter\textquotedblright{}.) Say it once more.
\end{tcolorbox}
\section[\texorpdfstring{\bn{টক} (ṭôk)}{টক (ṭôk)}]{\bn{টক} (ṭôk) — sour}

\label{lesson:BN-C75-tok}



\begin{quote}
Before the new one: say the Bengali for dirty, then the Bengali for bitter.
\end{quote}

\subsection*{You'll want to know: \bn{টক}}
\textbf{\bn{টক}} — \emph{ṭôk} — \textquotedblleft{}sour\textquotedblright{}.

\begin{cousinweb}
\textbf{\bn{টক}} (\emph{ṭôk}) is a tadbhava, a word inherited from Sanskrit and reshaped by sound change. It comes from \emph{takra}, \textquotedblleft{}buttermilk\textquotedblright{} (CDIAL 5617). CDIAL lists the Bengali word under a second sense it reconstructs, \textquotedblleft{}bitter\textquotedblright{}.
\end{cousinweb}

\begin{grammarlens}[title={What you've built}]
One more word for this chapter.
\end{grammarlens}

\subsection*{Your turn}
\begin{itemize}
\raggedright
  \item \emph{Say it:} \emph{ṭôk}
  \item \emph{Say it:} \emph{ṭôk}, once more
  \item \emph{Say it:} say \emph{ṭôk}
  \item \emph{Recall:} say the Bengali for dirty, then the Bengali for bitter, then say \emph{ṭôk} again
\end{itemize}

\begin{tcolorbox}[breakable,colback=teal!4,colframe=teal!35!black,title={Before you move on}]
What does \bn{টক} mean? (\textquotedblleft{}Sour\textquotedblright{}.) Say it once more.
\end{tcolorbox}
